\chapter{Lắng nghe và nói}
\markboth{Lắng nghe và nói}{Listening and Speaking}


\section{Thầy không cần em đồng ý; thầy cần em suy nghĩ.}
\begin{truyen}{Thầy không cần em đồng ý; thầy cần em suy nghĩ.}{The teacher does not need your agreement; the teacher needs your thinking.}
\chuhoa{C}{âu ấy không nhằm ép ai gật đầu. Nó chỉ làm cả lớp dừng lại một nhịp, để mỗi người tự hỏi mình đang hiểu vấn đề bằng suy nghĩ của chính mình hay chỉ lặp lại điều quen nghe.}
\textit{(The sentence was not meant to force anyone to agree. It simply made the whole class pause for a moment, so that each student could ask whether they truly understood the matter or were only repeating what they were used to hearing.)}
\end{truyen}

\begin{baihoc}
Điều quý nhất trong lớp học không phải sự đồng thanh, mà là sự suy nghĩ thật.
\textit{(The most valuable thing in class is not uniform agreement, but genuine thought.)}
\end{baihoc}

\begin{ghinhoanh}
\textbf{Short English to remember:}

Think for yourself before you agree.

\end{ghinhoanh}

\begin{tuhocnhanh}
1. Khi em đồng ý với một ý kiến, em thường đồng ý vì hiểu hay vì quen? \textit{(When you agree with an idea, do you agree because you understand it or because you are used to it?)}

2. Câu này nhắc em thay đổi điều gì trong cách học? \textit{(What does this sentence ask you to change in the way you study?)}
\end{tuhocnhanh}
\ngancach


\section{Một câu hỏi hay có thể thay đổi cả buổi học.}
\begin{truyen}{Một câu hỏi hay có thể thay đổi cả buổi học.}{One good question can change an entire lesson.}
\chuhoa{C}{âu hỏi vừa được viết lên bảng, căn phòng lập tức chậm lại. Không khí buổi học đổi khác, vì từ lúc đó các bạn không còn chỉ chờ đáp án, mà bắt đầu đi tìm cách hiểu.}
\textit{(As soon as the question was written on the board, the room slowed down. The atmosphere of the lesson changed, because from that moment on the students were no longer waiting only for an answer; they began looking for understanding.)}
\end{truyen}

\begin{baihoc}
Một câu hỏi mạnh có thể mở ra tư duy tốt hơn nhiều lời giải sẵn.
\textit{(A strong question can open better thinking than many ready-made explanations.)}
\end{baihoc}

\begin{ghinhoanh}
\textbf{Short English to remember:}

Good questions reshape the whole lesson.

\end{ghinhoanh}

\begin{tuhocnhanh}
1. Em từng nghe câu hỏi nào làm mình suy nghĩ lâu hơn bình thường? \textit{(Have you ever heard a question that made you think longer than usual?)}

2. Nếu được đặt một câu hỏi cho cả lớp, em sẽ muốn hỏi điều gì? \textit{(If you could ask one question to the whole class, what would you ask?)}
\end{tuhocnhanh}
\ngancach


\section{Im lặng không phải trống rỗng — im lặng là nơi ý tưởng ra đời.}
\begin{truyen}{Im lặng không phải trống rỗng — im lặng là nơi ý tưởng ra đời.}{Silence is not emptiness; it is where ideas are born.}
\chuhoa{C}{âu ấy đi qua lớp rất nhẹ. Chỉ có tiếng quạt và tiếng lật vở còn lại, nhưng chính trong khoảng yên đó, nhiều bạn bắt đầu nảy ra cách nghĩ mới mà lúc ồn ào không ai nhận ra.}
\textit{(The sentence passed through the classroom very gently. Only the fan and the turning pages remained, yet in that quiet space many students began to form new ideas that no one noticed while the room was noisy.)}
\end{truyen}

\begin{baihoc}
Không phải mọi khoảng lặng đều rỗng; có khoảng lặng đang làm việc cho trí óc.
\textit{(Not every silence is empty; some silences are working for the mind.)}
\end{baihoc}

\begin{ghinhoanh}
\textbf{Short English to remember:}

Ideas often begin in silence.

\end{ghinhoanh}

\begin{tuhocnhanh}
1. Em có thường sợ khoảng im lặng trong lớp không? Vì sao? \textit{(Do you often feel uncomfortable with silence in class? Why?)}

2. Khi cần nghĩ sâu hơn, em cần loại không gian nào? \textit{(What kind of space do you need when you want to think more deeply?)}
\end{tuhocnhanh}
\ngancach


\section{Em học được gì khi em im lặng lắng nghe?}
\begin{truyen}{Em học được gì khi em im lặng lắng nghe?}{What do you learn when you listen in silence?}
\chuhoa{C}{âu thầy nói xong làm vài bạn đang định chen lời chợt ngừng lại. Khi không vội đáp, người ta nghe rõ hơn cả điều người khác nói lẫn điều mình đang nghĩ.}
\textit{(What the teacher said made a few students who were about to interrupt suddenly stop. When people do not rush to answer, they hear more clearly both what others are saying and what they themselves are thinking.)}
\end{truyen}

\begin{baihoc}
Lắng nghe tốt không chỉ giúp hiểu người khác, mà còn giúp hiểu chính mình.
\textit{(Listening well helps you understand not only others, but also yourself.)}
\end{baihoc}

\begin{ghinhoanh}
\textbf{Short English to remember:}

Quiet listening teaches more than quick speaking.

\end{ghinhoanh}

\begin{tuhocnhanh}
1. Lần gần đây nhất em thực sự lắng nghe là khi nào? \textit{(When was the last time you truly listened?)}

2. Điều gì khiến em hay vội nói trước khi nghe hết? \textit{(What usually makes you speak before listening to the end?)}
\end{tuhocnhanh}
\ngancach


\section{Đừng sợ im lặng; sợ nói mà không nghĩ.}
\begin{truyen}{Đừng sợ im lặng; sợ nói mà không nghĩ.}{Do not fear silence; fear speaking without thinking.}
\chuhoa{C}{âu ấy làm mấy tiếng đáp lời vội vàng khựng lại. Không khí trong lớp bớt gấp gáp ngay lập tức, vì ai cũng nhận ra điều đáng ngại không phải là vài giây yên, mà là một câu nói thiếu suy nghĩ.}
\textit{(The sentence made a few rushed replies stop short. The atmosphere in class became less hurried at once, because everyone realized that what should be feared is not a few seconds of silence, but a sentence spoken without thought.)}
\end{truyen}

\begin{baihoc}
Nói ít nhưng chín vẫn tốt hơn nói nhanh mà rỗng.
\textit{(Speaking little but thoughtfully is better than speaking fast and empty.)}
\end{baihoc}

\begin{ghinhoanh}
\textbf{Short English to remember:}

Silence is safer than careless words.

\end{ghinhoanh}

\begin{tuhocnhanh}
1. Em đã từng hối hận vì nói quá nhanh chưa? \textit{(Have you ever regretted speaking too quickly?)}

2. Em có thể tập dừng lại bao lâu trước khi trả lời? \textit{(How long could you train yourself to pause before answering?)}
\end{tuhocnhanh}
\ngancach


\section{Người hiểu thường ít nói — vì họ biết lời có trọng lượng.}
\begin{truyen}{Người hiểu thường ít nói — vì họ biết lời có trọng lượng.}{People who understand often speak less, because they know words carry weight.}
\chuhoa{C}{âu ấy được nói ra khi vài bạn đang tranh nhau phát biểu. Sau đó lớp tự nhiên trầm xuống, như hiểu rằng lời nói có giá trị nhất không phải lúc nào cũng là lời nói nhiều nhất.}
\textit{(The sentence was spoken while a few students were competing to answer. After that, the room naturally became calmer, as if everyone understood that the most valuable words are not always the most frequent ones.)}
\end{truyen}

\begin{baihoc}
Người biết trọng lời sẽ biết trọng cả lúc nên im lặng.
\textit{(A person who respects words will also respect the moments when silence is better.)}
\end{baihoc}

\begin{ghinhoanh}
\textbf{Short English to remember:}

Weighty words are never careless words.

\end{ghinhoanh}

\begin{tuhocnhanh}
1. Em có đang dùng lời nói quá dễ dãi không? \textit{(Are you using words too carelessly?)}

2. Người em thấy nói ít mà sâu là ai? \textit{(Who is someone you know that speaks little but says meaningful things?)}
\end{tuhocnhanh}
\ngancach


\section{Lớp học đáng nhớ là lớp học có những khoảnh khắc im lặng đầy suy nghĩ.}
\begin{truyen}{Lớp học đáng nhớ là lớp học có những khoảnh khắc im lặng đầy...}{A memorable class is a class with moments of thoughtful silence.}
\chuhoa{C}{âu thầy nói xong không khiến lớp ồn lên, mà khiến mọi thứ lắng lại. Chính những giây như thế làm buổi học ở lại lâu trong trí nhớ, vì ai cũng có phần suy nghĩ riêng của mình trong đó.}
\textit{(What the teacher said did not make the room louder; it made everything settle. Moments like that stay in memory for a long time, because everyone has a piece of personal reflection inside them.)}
\end{truyen}

\begin{baihoc}
Chiều sâu của buổi học nhiều khi được tạo ra bởi những phút không ai tranh nói.
\textit{(The depth of a lesson is often created in the minutes when no one rushes to speak.)}
\end{baihoc}

\begin{ghinhoanh}
\textbf{Short English to remember:}

The best classes leave space for reflection.

\end{ghinhoanh}

\begin{tuhocnhanh}
1. Em nhớ nhất buổi học nào có một khoảng im lặng đặc biệt? \textit{(Which lesson do you remember most because of a special quiet moment?)}

2. Theo em, điều gì làm một lớp học đáng nhớ? \textit{(What makes a class memorable in your view?)}
\end{tuhocnhanh}
\ngancach


\section{Câu trả lời đúng có thể sai thời điểm; câu hỏi đúng luôn đúng lúc.}
\begin{truyen}{Câu trả lời đúng có thể sai thời điểm; câu hỏi đúng luôn đún...}{A correct answer can come at the wrong time; a good question is always timely.}
\chuhoa{C}{âu ấy làm cả lớp nhận ra rằng học không chỉ là nhặt cho đúng một đáp án. Có lúc điều cần nhất không phải trả lời ngay, mà là đặt ra đúng câu hỏi để buổi học đi đúng hướng.}
\textit{(The sentence made the whole class realize that learning is not just about grabbing the correct answer. Sometimes what matters most is not answering immediately, but asking the right question so the lesson can move in the right direction.)}
\end{truyen}

\begin{baihoc}
Một câu hỏi đúng lúc có thể cứu cả một cách nghĩ đang đi sai đường.
\textit{(A timely question can rescue an entire line of thought that is going in the wrong direction.)}
\end{baihoc}

\begin{ghinhoanh}
\textbf{Short English to remember:}

The right question matters at the right time.

\end{ghinhoanh}

\begin{tuhocnhanh}
1. Em có thường cố trả lời thật nhanh thay vì hiểu vấn đề không? \textit{(Do you often try to answer quickly instead of understanding the issue?)}

2. Câu hỏi nào em nên đặt ra thường xuyên hơn khi học? \textit{(What question should you ask more often while studying?)}
\end{tuhocnhanh}
\ngancach


\section{Suy nghĩ là thứ em làm khi em không vội trả lời.}
\begin{truyen}{Suy nghĩ là thứ em làm khi em không vội trả lời.}{Thinking is what you do when you do not rush to answer.}
\chuhoa{C}{âu ấy làm nhịp lớp học chậm lại đúng mức cần thiết. Không ai thấy khó chịu vì phải đợi, bởi chính khoảng đợi ấy cho người học thời gian biến phản xạ thành suy nghĩ.}
\textit{(The sentence slowed the rhythm of the class down just enough. No one minded having to wait, because that waiting time allowed learners to turn reaction into thought.)}
\end{truyen}

\begin{baihoc}
Không vội là điều kiện để hiểu sâu hơn.
\textit{(Not rushing is a condition for deeper understanding.)}
\end{baihoc}

\begin{ghinhoanh}
\textbf{Short English to remember:}

Thinking starts when hurry stops.

\end{ghinhoanh}

\begin{tuhocnhanh}
1. Việc gì ở lớp sẽ tốt hơn nếu em chậm lại một chút? \textit{(Which task in class would improve if you slowed down a little?)}

2. Em có đang nhầm giữa nhanh và giỏi không? \textit{(Are you confusing being fast with being good?)}
\end{tuhocnhanh}
\ngancach


\section{Thầy nói xong; giờ đến lượt em nghĩ.}
\begin{truyen}{Thầy nói xong; giờ đến lượt em nghĩ.}{The teacher has finished speaking; now it is your turn to think.}
\chuhoa{C}{âu ấy khép lại phần giảng nhưng mở ra phần quan trọng hơn. Từ đó, trách nhiệm không còn ở lời thầy nữa mà chuyển sang đầu óc của từng học sinh.}
\textit{(The sentence closed the explanation but opened something more important. From that moment on, the responsibility no longer belonged to the teacher's words, but to the mind of each student.)}
\end{truyen}

\begin{baihoc}
Nghe chỉ là bước đầu; nghĩ tiếp mới là phần học của người học.
\textit{(Listening is only the first step; continuing to think is the learner's true part of study.)}
\end{baihoc}

\begin{ghinhoanh}
\textbf{Short English to remember:}

After teaching comes thinking.

\end{ghinhoanh}

\begin{tuhocnhanh}
1. Sau khi nghe giảng, em thường làm gì để biến điều nghe được thành điều hiểu được? \textit{(After listening to a lesson, what do you usually do to turn what you heard into what you understand?)}

2. Nếu phải tự tổng kết một bài học, em sẽ bắt đầu từ đâu? \textit{(If you had to summarize a lesson for yourself, where would you begin?)}
\end{tuhocnhanh}
\ngancach
